\documentclass[oneside]{book}

\usepackage{sectsty}
\usepackage{graphicx}

\usepackage{listings}
\usepackage{xcolor}

\lstset{
  upquote=true,
  basicstyle=\fontencoding{T1}\ttfamily\small,
  frame=single,
  numbers=left,
  numberstyle=\small,
  stepnumber=1,
  numbersep=15pt,
  numberstyle=\ttfamily\small,
  xleftmargin=2em,
  framexleftmargin=0pt,
  framesep=8pt,
  showstringspaces=false,
  keepspaces=true,
  columns=fullflexible,
  breaklines=true,
  tabsize=4
}

\lstset{escapeinside={(*@}{@*)}}

\usepackage{fancyvrb}

\usepackage{minted}
\usepackage{hyperref}

\usepackage{lipsum}% http://ctan.org/pkg/lipsum
\usepackage{titletoc}% http://ctan.org/pkg/titletoc

\titlecontents*{chapter}% <section-type>
  [0pt]% <left>
  {}% <above-code>
  {\bfseries\chaptername\ \thecontentslabel\quad}% <numbered-entry-format>
  {}% <numberless-entry-format>
  {\bfseries\hfill\contentspage}% <filler-page-format>

\title{From Scratch:\\The Design and Implementation of a SAT Solver From Scratch}
\author{Oliver Lie}
\date{\today}

\begin{document}
\maketitle

\tableofcontents

\chapter*{Preface}

Hello! Thank you for taking me seriously enough to pick up this book and begin reading it.
Before we continue, I just want you to know that I am not an expert when it comes to things 
such as algorithms, data structures, general computer science, or even SAT solvers 
(what this book is about). Despite that, that's exactly why I wanted to write this. 
These things that I am not an expert in are my passion, and by doing this technical research, 
I hope to become a little less of ``not an expert'' in those fields.

The aim of this book isn't to improve or revolutionize SAT solving algorithms. 
Rather, it is to make the methods and reasoning more accessible to the lay undergraduate 
(I am open to the possibility that I myself may be the lay undergraduate, and you are 
simply smarter than me).

Throughout this book we will build our own modern SAT solver from scratch, 
complete with all the bells and whistles. We will begin with the most basic 
algorithm imaginable, brute force, and gradually work our way through DPLL, CDCL, and beyond.

Throughout this journey we will encounter thought experiments, exercises, and a
handful of experiments intended to solidify our understanding. Any external resources will 
be cited (hopefully correctly), and in keeping with the goal of accessibility:
\begin{itemize}
    \item Questions that are explicitly asked will eventually be answered.
    \item Proofs will be explained in informal language before becoming formal.
    \item Code will be provided and linked through GitHub (or another hosting platform).
    \item Revisions to this book will be made whenever I discover mistakes or better explanations.
\end{itemize}

There is one more thing that deserves acknowledgement: some pieces of code were AI-assisted.
To be completely transparent, ``AI-assisted'' means that instead of spending days repeatedly 
banging my head against a wall trying to understand a bug or some subtle algorithmic detail, 
I occasionally asked an AI to explain concepts or help debug an implementation. Every word of this book, 
however, is my own. AI did not generate any of the explanatory text you will read here, 
although it certainly helped produce some of the code that accompanies it.

There are practical reasons for this. SAT solvers are difficult to implement correctly, 
and introducing subtle bugs is remarkably easy. In fact, before deciding to write this book, 
I had already attempted to build a SAT solver once before. I never managed to get a working CDCL implementation.

That failure, however, taught me a tremendous amount. So in a sense, we are learning
SAT solving together (isn't that comforting?). My previous mistakes have shown me many 
of the pitfalls, and my hope is that by documenting this journey I can help you avoid making the same ones.
So without further ado, let's begin. By the time we reach the end, I hope I will have 
taught you more than I myself learned.

\chapter{Introduction}

SAT solvers, in my opinion, are one of the most underrated pieces of algorithmic 
technology of the modern day. They are used in both software and hardware verification 
(circuitry testing), product configuration, package management, computational biology, 
 particle physics, solving graph problems, and much much more.
 \footnote{\url{https://www.cs.utexas.edu/~isil/cs389L/lecture4-6up.pdf}}

In a more computer science-focused world, they have applications in 
NP-Completeness. Since (nearly) every problem has a reduction to a 
Boolean SAT problem, creating one really really good SAT solver lets us 
solve all those problems in a ``fast'' manner. So they are important everywhere
despite the vast majority of us never really thinking about their existence.
One note on what ``fast'' means: even if solving a SAT problem takes a small amount of time,
the algorithmic complexity isn't necessarily ``fast'' as there is no known 
polynomial algorithm for any NP-Complete problem 
(I am a firm believer that $P = NP$, and yes I know that makes me crazy).

Before we start writing out code, proofs, and the like, we need to make sure we're 
all on the same page on things such as input and ideas. Reading through this 
introduction is important so that you understand the basic notation used throughout this book.

One last thing to note: we will write out pseudocode rather than the code for one specific programming 
language. However, although the idea behind pseudocode is that it's programming 
language agnostic (it doesn't carry artifacts of any real programming language (usually)), 
this idea is poorly executed when it comes to formatting. My algorithms professor will 
probably scream at me for this, but I'm making the executive decision that we will
follow a general C-based syntax for readability purposes. That is, 
we will use braces \{\} for block scope, parentheses () for grouping, \verb|=| 
for assignment, \verb|==| for equality, and \verb|<=|, \verb|>=|, etc.\ for their 
respective uses. Everything that can be inferred visually will be used. Another thing to note is that we will 
pretend our arrays can dynamically resize, so we won't pre-allocate space, and we will also assume 1-based 
indexing for an easier map from variables to indices, as you will see later.

\section{What Do SAT Solvers Even Solve?}

``What do SAT solvers even solve,'' I hear you say? SAT solvers solve \textit{boolean satisfiability} problems. 
Another way to say that is they solve \textit{propositional logic} problems. 
There are many forms of propositional logic, and you may be familiar with it. Here are some basic examples:

\begin{itemize}
    \item $P \rightarrow Q$, an implication statement
    \item $(A \land B)$, an and statement
    \item $(A \lor B)$, an or statement
    \item $\neg A$, a negation/not statement
\end{itemize}

Of course these are very basic examples, and we haven't included every operator such as xor, 
nand, etc. As you know, we can rewrite implications to contain only and, or, and not operators,
which is what SAT solvers operate on. In other words, SAT solvers only operate on Boolean algebra equations.
\footnote{A Boolean algebra equation is one containing only and ($\land$), or ($\lor$), 
and not ($\neg$) operations (at a basic level).}

However, SAT solvers don't operate on just any Boolean algebra equation. No no no, they must be in a 
proper form, called CNF (more on that in the next section). CNF (Conjunctive Normal Form) equations are a 
``conjunction of disjunctions,'' i.e., clauses of ors, combined with ands.

Some examples of CNF are as follows:

\begin{itemize}
    \item $(A \lor B)$
    \item $(A \lor B) \land (\neg C)$
    \item $(A \lor B) \land (\neg C \lor D)$
\end{itemize}

Some examples that are \textbf{not} CNF:

\begin{itemize}
    \item $\neg(A \lor B)$, where there is a NOT in front of a clause
    \item $\neg(A \lor B) \land (C)$, where there is a NOT in front of a clause
    \item $(A \lor B \land C)$, where there is an AND within a clause
    \item $(A \land B) \lor (C)$, where it is in DNF form.
    \footnote{DNF (Disjunctive Normal Form) is a ``disjunction of conjunctions,'' 
    i.e., clauses of ands combined by ors. Solving these is trivially easy, 
    as we just need to find one clause where no element evaluates to false.}
\end{itemize}

\section{CNF Encoding}

Now that we know what our inputs look like, the next question (I hope you're
asking questions) is ``how do we encode that?'' Well good thing you're reading
this, because that's what we're here to talk about! In this section, we'll expand
our base of boolean algebra, including some definitions, transformations, and
finally good encoding.

\subsection*{Definitions We'll Use}

\begin{description}
    \item[Clause] A clause is a collection of variables, in which each variable is
    separated by an or ($\lor$), and each clause is separated by an and
    ($\land$).

    \item[Variable] A variable is an element of a clause, taking on a value of
    True, False, or Unassigned. Each variable can be negated ($\neg$) or not,
    affecting its evaluated value.

    \item[Literal] A literal is an evaluated variable, evaluating to one of True,
    False, or Unassigned. Since a literal is an evaluated variable, it cannot be
    negated. A literal whose corresponding variable is unassigned evaluates to
    Unassigned, regardless of whether or not the variable is negated.
\end{description}

A quick example of the difference of a literal value, using all three of our
definitions:

\[
(\neg A)
\]

is a clause of one variable ``A'', which is negated. If the variable
$A=False$, then its literal value is $True$, because

\[
\neg False \equiv True.
\]

In other words, $A$ was assigned the value False, and then it evaluated to
True. This will be very important in the future, so please take the time to
understand it well enough.

As discussed earlier, a CNF equation only has three operators: and ($\land$),
or ($\lor$), and not ($\neg$), but in propositional logic, there are many more
operators, and their clauses are not guaranteed to be in the format we want.
There are many algorithms to transform any propositional logic equation into
CNF, but frankly, that's beyond the scope of this book (it's also beyond my
scope). Instead, we will just look at how we transform each operator into an
equivalent format using just our three operators.

\subsection*{Implication ($\rightarrow$)}

\begin{center}
\begin{tabular}{cc|c|c|c}
$p$ & $q$ & $\neg p$ & $p\rightarrow q$ & $\neg p\lor q$\\
\hline
T & T & F & T & T\\
T & F & F & F & F\\
F & T & T & T & T\\
F & F & T & T & T
\end{tabular}
\end{center}

\subsection*{Xor ($\oplus$)}

\begin{center}
\begin{tabular}{cc|cc|c|cc|c}
$p$ & $q$
& $\neg p$
& $\neg q$
& $p\oplus q$
& $p\lor q$
& $\neg p\lor\neg q$
& $(p\lor q)\land(\neg p\lor\neg q)$\\
\hline
T&T&F&F&F&T&F&F\\
T&F&F&T&T&T&T&T\\
F&T&T&F&T&T&T&T\\
F&F&T&T&F&F&T&F
\end{tabular}
\end{center}

\subsection*{Biconditional ($\leftrightarrow$)}

\begin{center}
\begin{tabular}{cc|c|cc|c}
$p$ & $q$
& $p\leftrightarrow q$
& $p\rightarrow q$
& $q\rightarrow p$
& $(\neg p\lor q)\land(\neg q\lor p)$\\
\hline
T&T&T&T&T&T\\
T&F&F&F&T&F\\
F&T&F&T&F&F\\
F&F&T&T&T&T
\end{tabular}
\end{center}

We won't look at negations of operators such as nand, nor, xnor, or any
others. I'm not a boolean algebra expert when it comes to what operators
exist. But if you were given an equation with these extra operators, you could
easily apply the correct equivalency, then apply an algorithm that converts any
boolean algebra equation into one of CNF form.\footnote{I honestly may be
talking out of my ass here. I will do more research on this subject at a later
date.}

\subsection*{DIMACS CNF Format}

Now that we know what CNF looks like, we can now discuss how it's formatted
for a SAT solver's consumption. Typically, CNF is formatted in
\emph{DIMACS CNF} form.

Some rules for DIMACS CNF form are:

\begin{enumerate}
    \item The file may begin with comment lines. The first character of each
    comment line must be a lower case letter ``c''. Comment lines typically
    occur in one section at the beginning of the file, but are allowed to appear
    throughout the file.

    \item The comment lines are followed by the ``problem'' line. This begins
    with a lower case ``p'' followed by a space, followed by the problem type,
    which for CNF files is ``cnf'', followed by the number of variables followed
    by the number of clauses.

    \item The remainder of the file contains lines defining the clauses, one by
    one.

    \item A clause is defined by listing the index of each positive literal, and
    the negative index of each negative literal. Indices are 1-based, and for
    obvious reasons the index 0 is not allowed.

    \item The definition of a clause may extend beyond a single line of text.

    \item The definition of a clause is terminated by a final value of
    ``0''.

    \item The file terminates after the last clause is defined.
\end{enumerate}

There are some things about this syntax to be wary of:

\begin{enumerate}
    \item The definition of the next clause normally begins on a new line, but
    may follow, on the same line, the ``0'' that marks the end of the previous
    clause.

    \item In some examples of CNF files, the definition of the last clause is
    not terminated by a final ``0''.

    \item In some examples of CNF files, the rule that the variables are
    numbered from 1 to $N$ is not followed. The file might declare that there
    are 10 variables, for instance, but allow them to be numbered 2 through
    11.

    \item Comments can appear anywhere in the file, but the comment must be
    on its own line and it must start with ``c''.
    \footnote{\url{https://jix.github.io/varisat/manual/0.2.0/formats/dimacs.html}}
\end{enumerate}
For a quick example, here is a simple way to format a DIMACS CNF equation:

\begin{center}
  \begin{BVerbatim}
    p cnf 4 2
    1 2 -3 0
    -1 4 0
  \end{BVerbatim}
\end{center}

And one with comments:

\begin{center}
  \begin{BVerbatim}
    c this is a comment
    p cnf 4 2
    c this is another comment
    1 2 -3 0
    -1 4 0
  \end{BVerbatim}
\end{center}

In order to increase the compatibility of our SAT solver, we will be bending
some rules.

We will firstly allow missing header lines. In the case where no header is
passed in, we will instead infer the number of clauses and variables. We will
also be very lenient in where comments can appear, but we will enforce that
variables must be between $1$ and $N$.

\begin{enumerate}
    \item Comments can appear anywhere in the file, but they \textbf{MUST}
    be on their own line and start with ``c''.

    \item A header is not necessary, and if one is not provided, then the
    number of variables and clauses shall be inferred. However, if one is
    provided, then the number of variables and clauses \textbf{MUST} match the
    header.

    \item All clauses, with the exception of the last clause, must end with
    ``0'', and can extend multiple lines.

    \item All variables are numbered $1$ through $N$, where $N$ is the number
    of variables. That is, if you say there are 10 variables, they must be
    labeled $1,2,\ldots,10$ and not $2,3,\ldots,11$ or any other convention.

    \item (For compatibility) Upon encountering a ``\%'' character outside of
    a comment, we will stop parsing immediately.
\end{enumerate}

Thus, valid input can take on many forms:

\begin{center}
  \begin{BVerbatim}
    p cnf 4 2
    1 2 -3 0
    -1 4 0
  \end{BVerbatim}
\end{center}

\begin{center}
  \begin{BVerbatim}
    4 2
    1 2 -3 0
    -1 4 0
  \end{BVerbatim}
\end{center}

\begin{center}
  \begin{BVerbatim}
    p cnf 4 2
    c this comment doesn't have to be at the top
    1 2 -3 0
    -1 4 0
  \end{BVerbatim}
\end{center}

\begin{center}
  \begin{BVerbatim}
      p cnf 4 2
      c this comment doesn't have to be at the top
      1 2 -3 0
      -1 4
      c the first two numbers are still the number
      c of variables and the number of clauses
  \end{BVerbatim}
\end{center}

\begin{center}
  \begin{BVerbatim}
    p cnf 4 2 1 2 -3 0
    -1 4 0
  \end{BVerbatim}
\end{center}

More information can be found here.\footnote{\url{https://people.sc.fsu.edu/~jburkardt/data/cnf/cnf.html}}

\section{Reading a CNF Equation Programmatically}

Now that we know what our input looks like, we can finally discuss reading it
programmatically. Firstly we must discuss how our SAT solver will be represented
in code. At the moment, we just need to keep track of

\begin{itemize}
    \item the number of variables
    \item the number of clauses
    \item each clause in our equation
    \item whether or not we've tried to solve the equation yet
    \item what our solution is
\end{itemize}

We also want some methods to parse the equation, construct a SAT solver,
solve the equation, print out the solution if it is satisfiable, and whether or not
there's a solution. Therefore, our object will look like so:

\begin{lstlisting}
public interface:
    SATSolver(string input)
    bool solveCNF()

private interface:
    void parse(string equation)
    int numVars
    int numClauses
    int[][] clauses
    optional<bool>[] vars
\end{lstlisting}

Our pseudocode is as follows:

\begin{lstlisting}
SATSolver(string input)
{
    vars = no-value

    if (input is file)
    {
        string content = readFile(input)
        parse(content)
    }
    else
    {
        parse(input)
    }
}
\end{lstlisting}

\begin{lstlisting}
parse(equation)
{
    // The method of parsing if there's a header
    if (contains(equation, "p cnf"))
    {
        int[1..N] Arr

        for each (line in input)
        {
            if (startsWith(line, 'c'))
            {
                skip to next line
            }
            else
            {
                for each (token in line)
                {
                    append(Arr, token)
                }
            }
        }

        numVars = Arr[1]
        numClauses = Arr[2]

        idx = 2

        for (i = 1 up to numClauses)
        {
            int[1..N] clause

            while (Arr[idx] != 0)
            {
                append(clause, Arr[idx])
                idx = idx + 1
            }

            append(clauses, clause)
        }
    }

    // Method of parsing when there's no header
    else
    {
        numVars = 0
        numClauses = 0

        int[1..N] Arr

        for each (line in input)
        {
            if (startsWith(line, 'c'))
            {
                skip to next line
            }
            else
            {
                for each (token in line)
                {
                    numVars = max(numVars, token)

                    if (token == 0)
                    {
                        numClauses = numClauses + 1
                    }

                    append(Arr, token)
                }

                idx = 0

                for (i = 1 up to numClauses)
                {
                    int[1..N] clause

                    while (Arr[idx] != 0)
                    {
                        append(clause, token)
                        idx = idx + 1
                    }

                    append(clauses, clause)
                }
            }
        }
    }
}
\end{lstlisting}

This code will be implemented slightly differently in the example code, as I
am using C++. There is also a big bug, which is that the code assumes that
because ``p cnf'' appears in the file, there is a header. This is not always true
because it does not consider whether or not the header is on the same line as a
comment, so this will break it:

\begin{center} 
  \begin{BVerbatim}
      c p cnf
      1 2 -3 0
      -1 4 0
  \end{BVerbatim}
\end{center}

Because it assumes that there is a header when in fact there is not. I will
leave fixing this bug as an exercise to both the reader and the writer, however,
I will not be implementing this fix because if you're purposely trying to break
the constructor, you're not trying to solve a CNF equation anyways. The fix,
in case you're wondering, would be to just scan the lines and for each line that
doesn't start with ``c,'' check that it starts with ``p cnf,'' and if it does, then you
know for certain that there is a header. Now the real exercise is to determine
whether or not there's a header in just one pass instead of two.

\chapter{Brute Force}

Since the dawn of man, whenever encountered with a difficult problem, we have
decided ``eh, let's do it in the hardest way possible because I can't think of a
better solution.'' The history of brute forcing things goes all the way back to
our ancestors discovering fire,\footnote{This is probably not true.} all the way
to our current solution to climate change.\footnote{It's doing nothing and hoping
everything sorts itself out, which is really disappointing.}

The idea of brute forcing something is really as simple as it gets. We have $n$
number of literals, and therefore $2^n$ possible assignments, and we will check
them all until we have either found a solution, or until we've run out of
assignments to check. Because of how simple the idea is, it would just be
better to show code and explain that instead of explaining then showing code.

\begin{lstlisting}
solveCNF()
{
    bool[1..numVars] Assignment

    for (i = 0 up to 1 << numVars)
    {
        for (j = 0 up to numVars - 1)
        {
            Assignment[j] = i & (1 << j)
        }

        bool solved = true

        for each (clause in clauses)
        {
            bool satisfied_clause = false

            for (j = 1 up to clause.size)
            {
                int variable = clause[j]
                int v_idx = abs(variable)

                if (variable > 0)
                {
                    satisfied_clause = Assignment[v_idx] 
                      || satisfied_clause
                }
                else
                {
                    satisfied_clause = !Assignment[v_idx] 
                      || satisfied_clause
                }
            }

            solved &= satisfied_clause
        }

        if (solved == true)
        {
            vars = Assignment
            return true
        }
    }

    return false
}
\end{lstlisting}

Unfortunately, the formatting across the page kind of sucks, but let's try to
explain it anyways.

First, we create an array to hold each boolean value, and then next we start
our loop from $0$ all the way to ``1 << numVars.'' If you are unfamiliar with
what the left bit shift operator (`<<`) is, essentially, it is a fast way of
computing some number times $2$ times the number of places to shift. That is
all the explanation I am going to give because that's not the point. But by
computing `1 << numVars`, we are computing $1 \rightarrow 2^n$, which
happens to be the total number of possible variable combinations (from all
False to all True).

Then in the first inner for loop, we put the current variable combination into
the `Assignment` array. The index variable `i` holds the current variable
combination, and we isolate the $j$th variable using `1 << j` and put it into
the $j$th spot in the `Assignment` array.

Next, we evaluate the CNF equation against the current assignment. We first
assume that the equation is solved (innocent until proven guilty), and then we
evaluate each clause. We initially assume the clause is false because logically,
if one literal is True, then OR-ing it with False makes the ``clause satisfied''
condition true (we could exit out of the evaluation loop, but it's easier to read
if we don't). Then, we ``and'' the ``clause satisfied'' evaluation with the
``equation solved'' condition.

At the end of the evaluation loop, we check if we found a satisfiable
assignment. If we have, we then make our `vars` field, which holds the solution
if there is one, the assignment, and then we return True to indicate that we
have found a satisfiable assignment. If each possible assignment is not
satisfiable, then we return False, and `vars` will hold nothing.

Since it's trivially easy to see why the code works, we will move onto something
more interesting: optimizing brute force.

\section{Optimizing Brute Force}

I'll be completely honest with you, if you want to skip this section, go ahead.
It's more or less just something interesting enough we can do to make the
(arguably) worst algorithm\footnote{I do not know of any algorithm worse than
brute forcing.} better.

Since we're dealing with pseudocode, we're not going to consider the hardware
side of computers such as branch prediction, cache locality, or anything that
the hardware uses to execute our code (other than the CPU). We will consider
it from a pure Big-O perspective (excluding constant factors unless we're
choosing between two algorithms with the same computational complexity).

The first thing we will add is short circuiting to our clause evaluation. After
line 26 where we finish OR-ing the ``clause satisfied'' boolean, we will add in
a new block:

\begin{lstlisting}
...
bool satisfied_clause = false;

for (j = 1 up to clause.size)
{
    int variable = clause[j]
    int v_idx = abs(variable)

    if (variable > 0)
    {
        satisfied_clause =
            Assignment[v_idx] || satisfied_clause
    }
    else
    {
        satisfied_clause =
            !Assignment[v_idx] || satisfied_clause
    }

    if (satisfied_clause == True)
    {
        continue
    }
}

...
\end{lstlisting}

This just says ``once this clause is satisfied, go evaluate the next one
immediately.''

We can then add another conditional after this loop, which will move onto the
next variable assignment the instant the equation is no longer satisfied:

\begin{lstlisting}
...
bool solved = true;

for each (clause in clauses)
{
    bool satisfied_clause = false;

    for (j = 1 up to clause.size)
    {
        ...
    }

    if (satisfied_clause == True)
    {
        continue
    }

    solved &= satisfied_clause;

    if (solved == False)
    {
        break
    }
}

...
\end{lstlisting}

The next algorithm optimization we can make is removing the loop where we
move the current assignment into the `Assignment` array. Rather than doing it
upon every iteration, why don't we do it once we find a satisfiable assignment?

\begin{lstlisting}
solveCNF()
{
    bool[1..numVars] Assignment

    for (i = 0 up to 1 << numVars)
    {
        bool solved = true

        for each (clause in clauses)
        {
            bool satisfied_clause = false

            for (j = 1 up to clause.size)
            {
                int variable = clause[j]
                int v_idx = abs(variable)

                if (variable > 0)
                {
                    satisfied_clause = (i & (1 << j)) 
                      || satisfied_clause
                }
                else
                {
                    satisfied_clause = !(i & (1 << j))
                      || satisfied_clause
                }

                if (satisfied_clause == True)
                {
                    continue
                }
            }

            solved &= satisfied_clause

            if (solved == False)
            {
                break
            }
        }

        if (solved == True)
        {
            for (j = 0 up to numVars - 1)
            {
                Assignment[j] = i & (1 << j)
            }

            vars = Assignment
            return true
        }
    }

    return false
}
\end{lstlisting}

With that, that's all the optimization I can think of. In case you're
wondering, no this does not change the asymptotic runtime of the method,
however, it's easy to see why these changes would presumably make the
algorithm run faster as we're adding code that stops evaluation the moment
we've figured just enough out to know if it's worth continuing evaluation and
deferring extra calculation (recording the current assignment) until once we
know it's satisfiable.

As an exercise to the reader, I will ask you, how can you parallelize this
algorithm? And once you've figured that out, can you program it in a way
where once one thread finds a satisfactory assignment, it terminates all other
threads immediately instead of waiting for the others to finish?

\section{Brute Force -- Results}

To see how well our current algorithm does, running some tests would be a
good idea. Why run tests when we have already proven that our algorithm
works ``perfectly?'' To see how quickly it runs of course.

To test its speed (and as we add more features, to ensure correctness is
kept), we use a database of problems from here:

\begin{center}
\url{https://www.cs.ubc.ca/~hoos/SATLIB/benchm.html}
\end{center}

We gave our current implementation 1000 SAT instances (uf20-91), each
with 20 variables and 91 clauses each, all satisfiable. It took this baseline
brute force algorithm 2,119,301 milliseconds, which is roughly 35.32
minutes. This result is actually from the first time I wrote a brute force 
algorithm. The brute force result you can find on GitHub, it actually took 
1,770,851 milliseconds, which is about 29.51 minutes. I have no idea what I 
did in the rewrite that shaved off 4 minutes. I won't even lie to you reader,
I am a busy undergraduate student, so I didn't even do the "best practice" for 
testing, which would be to quit all other processes and then benchmark it. However,
while I was testing, I was doing a Canvas assignment, so I was doing lots of 
background work.

Normally in science you have to run an experiment at least three times for
data accuracy. But you must understand that I don't want to give up an hour
and a half of my life to an algorithm that we're going to be improving upon
anyways.

For fun, I will benchmark the short-circuit and parallel versions as well.
Short-Circuited (Basic): 49874ms (~50 seconds)

Parallel: 10589ms (~10.5 seconds)

For the parallel version, the implementation isn't perfect. I believe that 
we don't prevent false sharing amongst cache lines (ie we get more cache swaps between 
threads leading to more cache misses), which could explain why on my M1 MacBook Air, 
we don't get close to a ~8 times speedup. Since this was for fun, I'm not going to perfect 
the parallelized version, that can an exercise to the reader :)

\chapter{DPLL}

The Davis-Putnam-Logemann-Loveland (DPLL) algorithm is a complete
backtracking-based search algorithm for finding satisfiable assignments for SAT
problems. It is, dare I say, the foundational algorithm for modern SAT solvers,
as the state-of-the-art algorithm used by basically every SAT solver builds upon
the foundations that DPLL provides, as we will see later on.

It consists of three components:

\begin{itemize}
    \item Recursive backtracking
    \item Unit propagation
    \item Pure literal elimination
\end{itemize}

How does it work?

Imagine a binary tree of variables, each with two branches representing their
assignment, one for true and one for false. Suppose we don't know anything
about our CNF equation. That is, from looking at it, any variable could have
any value, and we don't know how each assignment helps us in finding out
whether or not the equation is satisfiable. How can we figure out if we can
satisfy the equation without just brute forcing it?

The answer is making an ass out of you and me to make assumptions.

\begin{figure}[H]
    \centering
    % TODO: Recreate DPLL decision tree figure.
    % Placeholder until the original figure is reconstructed.
    \fbox{\parbox[c][2.8in][c]{4.5in}{\centering
    DPLL Decision Tree}}
    \caption{Example DPLL decision tree (the outcomes may not correspond to an
    actual CNF instance, sue me).}
\end{figure}

Going back to our hypothetical, we don't know anything about anything. We
could have a giant decision tree to explore with an exponential number of leaves
to evaluate.

How does making assumptions help us?

Do we just run a random number generator and if it returns 1, we go ``okay,
it's satisfiable,'' and false otherwise?

No.

We make assumptions about our literal assignments.

Unlike brute forcing where we try every possible combination of variable
assignments at once, in DPLL we start with a literal, we'll call it $A$. There is
nothing in the equation that indicates whether or not $A$ should be true or
false; either could be valid.

So to tackle our problem, initially we will assume $A$ is true and see where
that leads us. We propagate that value throughout our clauses, and suppose
there is a clause that can't be satisfied. That's a contradiction, because if the
equation was satisfiable then $A = True$ couldn't have led us there.

So we've figured out that at the very least $A = True$ leads to a
contradiction.

We've now eliminated half of the search space immediately and only need to
consider the cases where $A = False$.

This may not make total sense immediately, so instead of giving one explanation
of the algorithm and then building it, we will build up the algorithm and explain
how each part works.

\footnote{Note for future me: write a better introduction.}

\footnote{\url{https://en.wikipedia.org/wiki/DPLL_algorithm}}

\section{Recursive Backtracking}

First, we need to rewrite our brute force algorithm to be recursive, as that's
the foundation for the entire algorithm (duh). To do this, we will introduce a
new method \textit{solve()} which will be our recursive solving method. To
also make things simpler, we will extract the equation evaluation into its own
separate method as well.

This means our SAT solver becomes like so:

\begin{lstlisting}
public interface:
    SATSolver(string input)
    int getNumVars()
    int getNumClauses()
    int[][] getClauses()
    optional<bool>[] getSolution()
    void printSolution()
    bool solveCNF()

private interface:
    void parse(string equation)
    bool solve(bool[] assignment)
    bool evaluate(bool[] assignment)
    int numVars
    int numClauses
    int[][] clauses
    optional<bool>[] lits
\end{lstlisting}

The flow essentially becomes like this: \textit{solveCNF()} calls
\textit{solve()}, which calls itself over and over and returns whatever solution
it finds back up to \textit{solveCNF()}, who will handle any sort of cleanup.

Why can't we just make \textit{solveCNF()} call itself recursively?

Good question.

That's because \textit{solveCNF()} will handle any setup that
\textit{solve()} will need in order to function properly, and it will handle any
cleanup that \textit{solve()} may create.

Let's now discuss how \textit{solve()} will work.

It will be passed in a \textit{bool[]} for the current assignment of literals.
When the size of our assignment is equal to the number of variables, we will
then evaluate the satisfiability of the assignment, and if it's satisfiable we
return true. Otherwise, we return false, backtrack, retry, until we either find a
satisfiable assignment or we've exhausted them all.

The pseudocode is as so, starting with \textit{evaluate()}:

\begin{lstlisting}
bool evaluate(bool[] assignment)
{
    bool solved = true;

    for each (clause in clauses)
    {
        bool satisfied_clause = false;

        for (j = 1 up to clause.size)
        {
            int variable = clause[j];
            int v_idx = abs(variable);

            if (variable > 0)
            {
                satisfied_clause = assignment[v_idx]
                  || satisfied_clause;
            }
            else
            {
                satisfied_clause = !assignment[v_idx] 
                  || satisfied_clause;
            }

            if (satisfied_clause == True)
            {
                continue;
            }
        }

        solved &= satisfied_clause;

        if (solved == False)
        {
            break;
        }
    }

    if (solved == True)
    {
        lits = assignment;
    }
    else
    {
        lits = no-value;
    }

    return solved;
}
\end{lstlisting}

As you can see, all we've done is move our equation evaluation from brute force
into its own function. One thing to note is that this method also handles
updating our \textit{lits} field. Thus, if \textit{evaluate()} returns true, then
\textit{lits} will hold the satisfying assignment, and if it returns false, then
\textit{lits} will hold nothing.

Now our recursive \textit{solve()} method:

\begin{lstlisting}
bool solve(bool[] assignment)
{
    if (assignment.size == numVars)
    {
        return evaluate(assignment);
    }

    append(assignment, True);

    if (solve(assignment))
    {
        return true;
    }

    assignment.back() = False;

    if (solve(assignment))
    {
        return true;
    }

    truncate(assignment, 1);
    return false;
}
\end{lstlisting}

This might be slightly confusing at first, so let's go over it.

We first check our base case. If our assignment has all of its variables, we
evaluate it and return the result whether it be true or false. And you'll notice
that whenever we return true from a call of \textit{solve()}, we immediately
return true afterwards, so therefore, if we find a satisfiable assignment, we
will go up our call stack back to \textit{solveCNF()}, so no need to worry that
it will terminate if we find a satisfiable assignment.

However, if our assignment still needs some variables, we append true to our
assignment, then we do a recursive call on \textit{solve()}. If that recursive
call leads to a satisfiable assignment, we return true, otherwise, we swap out
our last value for false and try \textit{solve()} again. And if that ends up
working out, we return true, but if it doesn't then we remove the last value
and return false.

So how do we know that it will guarantee termination for an unsatisfiable
problem?

Notice how when assigning the current variable true or false and it doesn't
work, we just remove the last one and go up the call stack by one? Well
suppose we're at the top of the call stack (we've assigned the first variable)
and the first variable is assigned false, and that didn't lead to a satisfiable
assignment. Then we remove that first variable from the assignment, and then
we return false. So our assignment array is empty, and we've returned to
\textit{solveCNF()}. Thus, \textit{solve()} will terminate eventually.

And how do we know that we'll try every single assignment?

Well for the very last variable, it's easy to see, if it being true doesn't satisfy
the equation, we assign it false, and if that doesn't work, we remove it and go
up the call stack. Then for the previous variable, if it was true, we assign it
false and go back to the last variable where it tries true and/or false again.
Then that doesn't work, we go up the call stack twice, and that continues until
we've found a satisfiable assignment or we've exhausted them all.\footnote{
This is a very informal proof, but what's most important is that we build up
our understanding of what's going on instead of dumping notation. Perhaps I
will do a formal proof in the future.}

Lastly, our \textit{solveCNF()} method:

\begin{lstlisting}
bool solveCNF()
{
    bool[1..numVars] assignment;

    return solve(assignment);
}
\end{lstlisting}

Because we're dealing with pseudocode, we don't have to do things such as
reserving or resizing the assignment array. And since \textit{evaluate()} takes
care of making sure our \textit{lits} field holds a valid assignment (if found)
and holds nothing if not, we don't need to do anything else in
\textit{solveCNF()}.

You may wonder though, why not pass in \textit{lits} to \textit{solve()} or
have no setup and make \textit{solveCNF()} purely recursive?

Those are very valid questions to ask.

This has to do with the differences between pseudocode and actual
implementation. In the actual implementation, some methods depend on the
state of \textit{lits}, such as whether or not it holds anything, and by passing
it into \textit{solve()}, we force it to hold something, even if there is no
satisfiable assignment. And it's true, that in \textit{solveCNF()} if we had no
setup, we could make it recursive on \textit{lits}, however, I chose to
implement it differently.

When you are reading this, please ask yourself these kinds of questions,
because oftentimes, the way you read my implementation is not the best way,
and even if it is, there is no reason as to why you can't think of different
implementations, try them out yourself, and compare. Perhaps, and I expect
this will happen often, you will find a better implementation than the one
written down here.

\subsection{Recursive Backtracking -- Results}

Running this code on the same 1000 instances (uf20-91), it took
1,846,644ms, which is roughly 30.78 minutes. This is actually a pretty big
improvement from our 35.32 minutes via brute forcing. Although this seems
like a big gain in speed, it should be noted that at this point the algorithm is
still just a fancier brute forcing algorithm. The next time I rewrote this algorithm,
it took 1,540,613ms, which is about 25.68 minutes. So once again, that's about a 
~5 minute speedup just from rewriting it. Perhaps we should rewrite every piece of 
software in the world to make it run 5 minutes faster. 

Here's a question for you, why might assigning our variables true instead of
false first lead to an increase in speed?

Hint: Look at the number of negations present in the testing set.

\section{Unit Propagation}

Unit propagation is the next step in implementing DPLL, and it introduces the
idea of implication.

Imagine clause: $A \rightarrow Z$, and we know that $A = False$. What
can we say about $Z$?

Well, if our equation is satisfiable, then the clause must be satisfiable,
therefore $Z = True$ because if it didn't, that would be a contradiction.

So instead of assigning variables $B$, $C$, \ldots, $X$, then assigning $Z$,
why don't we just do it now, and use that knowledge in other clauses
containing $Z$ or $\neg Z$?

That is essentially the idea behind unit propagation: the decisions we make
imply other assignments, and so we should force them now instead of waiting
until later.

However, this introduces a light hiccup into our implementation.

You see, our initial invariant about assigning variables is that they are
assigned in numerical order instead of (possibly) out of order.

``Not a problem! We'll just make the assignment array have space for every
variable so we can index it anywhere,'' I hear you think.

Unfortunately that interferes with our base case because it relies on the size
of the assignment array. If it is initially at a fixed size of
\textit{numVars}, then our base case will always fire because our assignment
array always has a size of \textit{numVars}. So we would either need to
change our base case or change how we store learned variables.

What about a learned set? One that holds whether or not we've learned a
variable (positive or negative), and then as we traverse down the call stack,
we just check if a variable is already learned, and if so, we give it its learned
value?

That would 100\% work.

In fact, the first time I implemented DPLL, that is exactly what I used.

However, it has the disadvantage of being annoying to maintain.

Instead, why don't we create a new data type called something like
\textit{var} that represents the state of each variable: \textit{True},
\textit{False}, and \textit{Unassigned}.

Then we can just maintain a simple counter of how many are not
\textit{Unassigned}, and then that keeps our base case nice and neat.

Yes dear reader, we are going to do that, however, if you'd like to use the
learned set, reading about how we do it without a learned set will help you
and offer clues on how to implement it your way.

\subsection{Changes To Our Implementation}

\begin{lstlisting}
    public interface:
        SATSolver(string input)
        int getNumVars()
        int getNumClauses()
        int[][] getClauses()
        optional<bool>[] getSolution()
        optional<var>[] getSolution()
        void printSolution()
        bool solveCNF()

    private interface:
        enum var
        {
            False
            True
            Unassigned
        }

        void parse(string equation)
        bool solve(bool[] assignment)
        bool solve(var[] assignment, int level)
        bool evaluate(bool[] assignment)
        bool evaluate(var[] assignment)
        void propagate()

        int numVars
        int numClauses
        int numAssigned

        int[][] clauses

        optional<bool>[] lits
        optional<var>[] vars
\end{lstlisting}

Here we just add in a new enum for the states that a variable can be, and
adjust the methods using \textit{bool}s to use the new enum type instead.

\subsection{``Textbook" Unit Propagation}

Before we start changing our algorithm, let's implement the
\textit{propagate()} method.

As stated before, the idea is that given a clause with either only one variable
or only one \textit{Unassigned} variable (and the rest evaluate to
\textit{False}), we can immediately say that \textit{Unassigned} variable will
be whatever value it needs to be to evaluate to \textit{True} for that clause.

In the ``textbook" way of implementation,
\footnote{\url{https://en.wikipedia.org/wiki/Unit_propagation}} 
once a clause is satisfied because of a literal, you remove every other
clause that is satisfied by that literal. And for every other clause containing 
the negation of the literal, you remove that negation entirely. This is done to shrink the 
search space.

We will implement this algorithm with as much malicious compliance as you can 
carry at an office job without getting fired:

\begin{lstlisting}
    bool solveCNF() {
        var[1..numVars] assignment

        for each (var in assignment) {
            var = Unassigned
        }

        propagate(assignment)
        solve(assignment, 0)
    }
\end{lstlisting}

We want to propagate before we start solving in case there are clauses with 
only one variable (it's already unit).

\begin{lstlisting}
    bool solve(bool[1..numVars] Assignment, int level) {
        if (level == numVars - 1) {
            return evaluate(Assignment)
        }

        bool[1..N] backupAssignment = assignment
        int[][] backupClauses = clauses

        if (Assignment[level] == Unassigned) {
            Assignment[level] = True
            propagate(Assignment)
        }

        if (solve(Assignment, level + 1)) {
            return True
        }

        clauses = backupClauses
        Assignment = backupAssignment

        if (Assignment[level] == Unassigned) {
            Assignment[level] = False
            propagate(Assignment)
        }

        if (solve(Assignment, level + 1)) {
            return True
        }

        return False
    }
\end{lstlisting}

Before we get onto the code for unit propagation, let's understand this code first. 
Because unit propagation no longer guarantees that variables are assigned in order, 
we need to pass in a ``level'' as a parameter so we know which variable we need to 
assign. We also now need to check that propagation didn't already assign this variable. 
If it's currently Unassigned, we need to assign it then propagate, but if it's been 
assigned before we assign it, then we know that it's the result of propagation, so we 
can continue onto the next level. Other than that, it's practically the same as our recursive 
backtracking code, except now we need to have a ``backup'' of our clauses and our assignment. 
Why? Because \textit{propagate()} \textbf{mutates} the clauses and the assignment in a way 
that breaks the assumption that our literals are assigned in ``numeric'' order. Since it 
returns nothing, we don't have a way to undo that easily, so instead, for readability, 
we just have a backup copy and whenever propagation leads to an unsatisfying assignment, 
we return false and undo whatever changes we made by copying back into our solver's fields.
And the reason why we don't need to undo the propagation before returning False is because 
when we move back up the call stack by one level, the method continues by undoing the effects for us.
Now let's look at how we perform textbook unit propagation:

\begin{lstlisting}
    void propagate(int[1..N] Assignment) {
        while (true) {
            Set<int> unitClauses;
            Set<int> unitVars;

            for (i = 0 up to numClauses) {
                clause = clauses[i]

                bool isUnit = True
                int unassignedIdx = -1
                int numUnassigned = 0

                for (j = 0 up to clause.size) {
                    int var = clause[j]
                    int idx = abs(var)

                    if (Assignment[idx] == Unassigned) {
                        numUnassigned = numUnassigned + 1
                        unassignedIdx = j
                    } else if (var > 0 && clause[j] == True || var < 0 && clause[j] == False) {
                        isUnit = false
                    }
                }

                isUnit = isUnit && numUnassigned == 1

                if (isUnit) {
                    add(unitClauses, i)
                    add(unitVars, clause[unassignedIdx])

                    int var = clause[unassignedIdx]
                    int idx = abs(var)

                    if (var > 0) {
                        assignment[idx] = True
                    } else {
                        assignment[idx] = False
                    }
                }
            }

            int clauseIdx = 0;
            for each (clause in clauses) {
                for each (var in clause) {
                    if (contains(unitVar, -var)) {
                        remove(clause, -var)
                    }

                    if (contains(unitVar, var) && contains(unitClause, clauseIdx)) {
                        remove(clauses, clause)
                    }
                }

                clauseIdx = clauseIdx + 1
            }

            if (unitClauses.size == 0 && unitVars.size == 0) {
                break
            }
        }
    }
\end{lstlisting}

Okay, now let's understand what's happening here. We immediately go into an 
infinite loop (nice), and we start looking at each clause for two things:
\begin{enumerate}
    \item The clause has only one Unassigned variable 
    \item Every other literal in the clause is False (remember, literals are evaluated variables)
\end{enumerate}

If both things are true, then the Unassigned variable will be assigned whatever value 
is required to make it evaluate to True within that clause, and we know we have found 
a unit variable and a unit clause. After we have found all of our unit variables and unit 
clauses, we then rescan every other clause and delete the each unit variable's negative,
(if we learn that -5 is unit, then we remove +5), and if the clause contains a unit 
variable, we just delete that clause entirely from the database. Once we stop finding 
new unit variables and new unit clauses, we stop entirely. You may wonder (I reread my 
code and thought I found a bug), ``how do we know that we always terminate?'' That's a 
great question. You see, this confusion may come from the fact that we rescan the database 
every iteration of the \textit{while()} loop searching for unit clauses. What is stopping 
us from re-finding the same clause as unit upon a new iteration? The answer is much simpler 
than you'd expect: since we assign our unit variables, the clauses upon the next iteration 
no longer foloow our criteria for being unit propagate-able because there is now a literal 
in that clause that is True. Therefore, the number of unit variables and clauses monotonically
decrease between iterations (a fancy way of saying the number can only go down).

Now, you may be wondering, ``what's the performance of our solver now?" Let me answer that 
as plainly as possible: worse than brute force. How? Easy! We are doing SO MANY things that 
hurt the computer. We are allocating TONS of memory in order to backtrack, we are allocating 
even more memory to propagate (hash sets aren't cheap), and hashing, while O(1), is not 
a cheap operation to do. So yeah, I can believe that it's worse than brute force because 
at least the brute force algorithm was easy on the CPU. Let's learn how this is 
\textit{usually}
\footnote{I wanted to say ``actually implemented'', however, that's a big generalization, and 
again, since I'm not the expert here, I don't think I have the authority to say that. The next 
section is how \textit{I} implemented it, and it roughly follows the idea on how SAT solvers 
implement this.}
implemented!

\section{Better Unit Propagation}

So the textbook style unit propagation sucked. Let's change that. As discussed before, 
the biggest bottleneck/inefficiency in the actual implementation is the number of copies 
we need to make of our clauses and assignments on each decision level. Assigning one variable 
creates a new decision level, so 20 variables means 20 levels. For an unsatisfiable 
CNF equation, I estimate we need to make about 3 copies of our clauses and assignments:
one for the initial storage, another for the first failure, and a third for the second. 
Even if we ignore the inefficiencies in how we propagate (memory allocations, traversals, and 
adjustments our arrays would need to make when we delete things), this is a huge amount of 
memory to store. If you think that 20 variables and 91 clauses is a lot, modern SAT solvers 
can handle equations with \textbf{millions} of variables, so creating an algorithm that 
makes brute force look attractive is actually an achievement.
\footnote{They might revoke my degree for that actually}

So what's the strategy? Firstly, it would be nice if we didn't actually mutate our 
clauses at all. Ie, no clause deletion, and no variable deletion. Next, it would be 
nice if we didn't need hash sets in order to store our unit variables. Instead, it
would be nice if we had a data structure where we put in our decided variables, and then 
it slowly shrinks as we iterate through. Ie, we propagate one on variable at a time and 
never again. This sounds like we need a \textit{queue}.
\footnote{That wasn't shoehorned in at all!}
Lastly, instead of copying our assignment, and constantly reiterating through the 
\textit{entire} clause database, what if we just marked which variables we assigned,
and which clauses we've satisfied? That way when we need to undo, we just mark those 
variables as Unassigned, and those clauses as no longer satisfied. That way we save a lot 
of memory (and besides, copying is NOT O(1) as you might be led to believe by that deceptively
simply ``='' operator).

Let's change our interface:

\begin{lstlisting}
    public interface:
        SATSolver(string input)
        int getNumVars()
        int getNumClauses()
        int[][] getClauses()
        optional<bool>[] getSolution()
        optional<var>[] getSolution()
        void printSolution()
        bool solveCNF()

    private interface:
        enum var
        {
            False
            True
            Unassigned
        }

        void parse(string equation)
        bool solve(var[] assignment, int level)
        bool evaluate(var[] assignment)
        (int[] vars, int[] clauses) propagate()

        int numVars
        int numClauses
        int numAssigned

        int[][] clauses
        bool[] satisfied

        optional<bool>[] lits
        optional<var>[] vars
\end{lstlisting}

We need to make some minimal changes to \textit{solveCNF()}:

\begin{lstlisting}
    bool solveCNF() {
        var[1..numVars] assignment

        for each (var in assignment) {
            var = Unassigned
        }

        propagate(assignment, 0)
        solve(assignment, 0)
    }
\end{lstlisting}

Since \textit{propagate()} now takes a decision level/variable index, we need to pass 
one to the method. Although no variable has an index of 0 (because it's reserved to mark the 
end of a clause), passing in $0$ just tells it ``Hey, there isn't a literal to propagate on,
so just find all clauses with only one variables and mark them as satisfied and `learn' those 
variables.'' And since this is at the root level, we don't need to undo anything if it turns out 
that the equation is unsatisfiable, since we're gathering things that HAVE to be a certain value 
if this equation is to be satisfiable.

Since we don't need to change \textit{solveCNF()}, we can go straight to \textit{solve()}:

\begin{lstlisting}
    bool solve(int[1..N] Assignment, int level) {
        if (level == numVars - 1) {
            return evaluate(Assignment);
        }

        if (Assignment[level] != Unassigned) {
            return solve(Assignment, level + 1)
        }

        Assignment[level] = True
        (int[] unitVars, int[] satisfiedClauses) = propagate(Assignment, level)

        if (solve(Assignment, level + 1) == True) {
            return True
        }

        for each (var in unitVars) {
            int idx = abs(var)
            Assignment[idx] = Unassigned
        }

        for each (clauseIdx in satisfiedClauses) {
            satisfied[clauseIdx] = False
        }

        Assignment[level] = False
        (int new_unitVars, int new_satisfiedClauses) = propagate(Assignment, level)

         if (solve(Assignment, level + 1) == True) {
            return True
        }

        for each (var in new_unitVars) {
            int idx = abs(var)
            Assignment[idx] = Unassigned
        }

        for each (clauseIdx in new_satisfiedClauses) {
            satisfied[clauseIdx] = False
        }

        Assignment[level] = Unassigned
        return False
    }
\end{lstlisting}

This method is practically the same, except now, we just check if we've already assigned 
the variable at the current decision level. If we have assigned a variable before we reach 
its decision level, then it must have been inferred/learned from the \textit{propagate()} 
method, so we don't need to call it again. Then, it's business as usual: assume the 
current variable is True, propagate, if it leads to a satisfiable assignment, return True, 
otherwise, we undo what we've assigned and marked as learned, then repeat with False. 

What's more interesting is our shiny new \textit{propagate()} method:

\begin{lstlisting}
    (int[] vars, int[] clauses) propagate(int[] assignment, int level) {
        queue<int> unit

        int[] unitVars
        int[] satisfiedClauses 

        unit.enqueue(level)

        while (unit.size > 0) {
            unit.dequeue()

            for (i = 0 up to numClauses) {
                if (satisfied[i]) {
                    continue
                }

                clause = clauses[i]

                bool isUnit = True
                int unassignedIdx = -1
                int numUnassigned = 0

                for (j = 0 up to clause.size) {
                    int var = clause[j]
                    int idx = abs(var)

                    if (Assignment[idx] == Unassigned) {
                        numUnassigned = numUnassigned + 1
                        unassignedIdx = j
                    } else if (var > 0 && clause[j] == True || var < 0 && clause[j] == False) {
                        isUnit = false

                        satisfied[i] = true
                        append(satisfiedClauses, i)
                    }
                }

                isUnit = isUnit && numUnassigned == 1

                if (isUnit) {
                    add(unitClauses, i)
                    add(unitVars, clause[unassignedIdx])

                    int var = clause[unassignedIdx]
                    int idx = abs(var)

                    if (var > 0) {
                        assignment[idx] = True
                    } else {
                        assignment[idx] = False
                    }

                    unit.enqueue(var)

                    append(unitVars, var)
                    satisfied[i] = true
                    append(satisfiedClauses, i)
                }
            }
        }

        return (unitVars, satisfiedClauses)
    }
\end{lstlisting}

So, let's discuss what's going on. \textit{propagate()} now takes in the decision level, 
which is also the same thing as an index for a variable (conceptually speaking). It then 
adds that to a queue, then while the queue is not empty, it scans the clauses as before. 
First it makes sure we're not operating on a clause that's already been satisfied 
(because we can't tell anything from it), and when it finds a clause that has a true literal, 
it marks the clause as satisfied. Then when it finds a unit clause, it marks the clause as 
satisfied, keeps track of the new unit variable its found, and then adds that unit variable 
to the queue. Then once it's done, it returns the unit variables and the indices of the satisfied 
clauses.

You may have some questions about this approach. Firstly, how do we make sure we don't 
ever add the same variable twice? The answer is the same as before: because we assign a 
variable once it's found to be in a unit clause, it can't be found as Unassigned again, 
and therefore can't be marked as a unit variable again. Then, since we skip over already 
satisfied clauses, the amount of clauses we scan each iteration decreases monotonically. 

And going back to our \textit{solve()} code, even if we mark the same clause as satisfied 
or the same variable as unit twice, it doesn't matter because when we backtrack, we simply 
restore it back to its state of unsatisfied or Unassigned, respectively. What's important is 
that we don't mark the same clause or variable twice \textit{across decision levels} because 
then, and only then, would our backtracking code become a problem because we wouldn't be 
correctly restoring the state of the equation back to what it was before the decision level 
was entered. But since we can't mark the same variable or clause twice \textbf{until}
we backtrack, there's nothing to worry about.

Before we end this section, you may also be wondering: can we use the \textit{satisfied}
array in our \textit{evaluate()} method to skip clauses we know we've already satisfied?
And the answer is yes, and I would encourage you to do so in your implementations.

\subsection{Better Unit Propagation - Results}
So after implementing better unit propagation, this took 868ms to solve all 1000 
satisfiable equations. The first time I implemented this version of unit propagation,
it actually took 2363ms, so this is somehow over twice as fast. And because it's so 
fast, just like the first time I implemented this, we will be now testing our solver on 
unsatisfiable equations as well.

Our second testing set comes from the same database as before.
\footnote{\url{https://www.cs.ubc.ca/~hoos/SATLIB/benchm.html}, in case you needed a reminder.}
And our next testing set is called "UUF50.218.1000." In other words, it's full of equations that 
have 50 variables, 218 clauses, and there are 1000 such equations in the set. In order to solve 
all 1000 equations, it took 140,411ms (~2.34 minutes). That's also much faster than the first time 
I implemented this algorithm. The first time I did, it took 646,228ms (~10.77 minutes). Honestly,
with this much speedup, I don't know if I was just bad at implementing these algorithms the first time,
\footnote{Highly likely}
or if downgrading my operating system (and subsequently wiping my harddrive clean)
\footnote{I HATED MacOS Tahoe. Sequoia is where it's at.}
just has that much of an effect on execution speed.
\footnote{Also likely, depending on how much RAM and swap is being used to begin with.}

Why didn't we test our previous versions out on unsatisfiable equations? That's a good question.
For brute forcers, since there is no satisfiable assignment, it needs to check every single 
one to conclude that it's unsatisfiable, which would take a lot of time. Likewise with our 
recursive backtracking solver. And since our (my) implementation of textbook unit propagation 
was \textit{worse} than brute force,
\footnote{A difficult feat.}
obviously we wouldn't want to test it on something that will take more time than a satisfiable 
equation. Also because the asymptotic runtime is \textit{still} exponential, having a much 
larger equation would take a much, much larger amount of time, even if they were satisfiable.

\section{Pure Literal Elimination}
\section{Optimizing Our Current Approach}
\section{Iterative DPLL}

\end{document}
